\chapter{Index Rebalancing}\label{ch:s1:index-rebalancing}

When a stock joins an index, every fund that tracks the index must buy it at the same close. Shleifer found in 1986 that stocks newly added to the S\&P 500
earned a significant abnormal return at the announcement, related to how much \omterm{def:s1:index-rebalancing:fund}{index funds} bought: demand curves for stocks slope down. Thirty years later
Patel and Welch found the effect reverting in the 2000s: added and deleted stocks no longer see a permanent shift in demand. On the synthetic market, \omterm{def:s1:index-rebalancing:fund}{index
funds} tracking 15\% of each member's shares push an addition by 5.3\% on average; a trader who buys at the announcement keeps 2.4\% of it when predictors
have taken half the move before the announcement, and 0.8\% when they have taken four fifths. This chapter is about index rules, predicting what they will
do, and trading an event that everyone sees coming. The build is \texttt{firm.indexevent}.

\section{Index rules and the calendar}

\begin{definition}[Index fund, rank day]\label{def:s1:index-rebalancing:fund}
An \emph{index fund}\index{index fund} holds the securities of an index in the index's weights, so that its trades follow the index's changes rather than
any view of the securities. A \emph{rank day}\index{rank day} is the date on which a rules-based index measures the characteristics (capitalisation, float,
liquidity) that decide its membership at the next reconstitution.
\end{definition}

Index changes (Book~1, chapter~15) come in two kinds. Indices whose changes are decided case by case, such as the S\&P 500 in the studies below, announce each change before it takes
effect.
Rules-based indices such as the Russell US indexes follow a published calendar, and anyone with the data can compute their changes.

\begin{dated}{2026-09}{The Russell reconstitution calendar}\label{dat:s1:index-rebalancing:russell}
FTSE Russell reinstated a semi-annual reconstitution of the Russell US indexes in 2026. The June 2026 reconstitution used April 30 as its \omterm{def:s1:index-rebalancing:fund}{rank day},
published preliminary lists on May 22 and took effect after the close on Friday, June 26. The December 2026 reconstitution has its \omterm{def:s1:index-rebalancing:fund}{rank day} on Friday,
October 30, preliminary additions and deletions after the close on November 13, a lockdown from November 30, and takes effect after the close on Friday,
December 11. FTSE Russell reports about \$12.2 trillion of investor assets benchmarked to or invested in products based on the Russell US indexes.
\end{dated}

The synthetic index follows the same pattern. It holds the 300 largest of the synthetic market's stocks, is rebuilt every 126 trading days (seventeen times
over the sample), and bands membership: a non-member joins only if it ranks 285th or better on the \omterm{def:s1:index-rebalancing:fund}{rank day}, and a member leaves only if it ranks worse
than 315th. The band (the buffer rule of Book~1, chapter~15) keeps stocks near the cut-off from churning in and out. Preliminary lists come sixteen trading
days after the \omterm{def:s1:index-rebalancing:fund}{rank day}, and the changes take effect at the close forty trading days after it: 331 additions and 270 deletions in all.

\section{Predicting additions and deletions}

\begin{definition}[Membership prediction]\label{def:s1:index-rebalancing:prediction}
\emph{Membership prediction}\index{membership prediction} estimates, before the \omterm{def:s1:index-rebalancing:fund}{rank day}, the probability that each stock will be in the index after the
reconstitution, by applying the index's rules to the stock's characteristics projected forward.
\end{definition}

On the \omterm{def:s1:index-rebalancing:fund}{rank day} itself a rules-based change is known exactly: applying the rule to that day's capitalisations reproduces every addition and deletion. Before
it, the capitalisations still have to move. \texttt{firm.indexevent.predict} simulates each stock's capitalisation forward with its own volatility, two
hundred times, and applies the rule to each draw; a stock is predicted to join if it joins in more than half the draws.

\begin{center}
\small
\begin{tabular}{@{}lcccc@{}}
\toprule
predicted before the \omterm{def:s1:index-rebalancing:fund}{rank day} & additions: precision & recall & deletions: precision & recall\\
\midrule
on the \omterm{def:s1:index-rebalancing:fund}{rank day} & 100\% & 100\% & 100\% & 100\%\\
20 trading days before & 86\% & 72\% & 86\% & 74\%\\
60 trading days before & 76\% & 45\% & 76\% & 46\%\\
\bottomrule
\end{tabular}
\end{center}

Twenty days out, 86\% of the predicted additions happen and 72\% of the actual additions were predicted; sixty days out, three quarters of the predictions
still come true but fewer than half the additions are found. The trade-off is the chapter's: the earlier the prediction, the less certain it is, and the
more of the price move is still ahead of it.

\section{Trading the event and the reversal}

The synthetic event is planted on top of each stock's market-adjusted return. \omterm{def:s1:index-rebalancing:fund}{Index funds} track a share of each member's shares and trade the changes at the
effective close; their demand's price effect is the square-root law of Book~7, chapter~27, $0.7\,\sigma\sqrt{Q/V}$, with $Q$ the shares they must buy and
$V$ the stock's daily volume. A share of the effect (\texttt{early}) is taken before the announcement by traders who predicted the change; the rest accrues
between the announcement and the effective close; half of it reverses over the twenty days after. With 15\% of shares tracked, the average push is 5.3\%.

\begin{omfigure}
\begin{tikzpicture}
\begin{axis}[width=11cm,height=5.4cm,xlabel={\small trading days from the announcement},ylabel={\small cumulative abnormal return (\%)},
  tick label style={font=\footnotesize},xmin=-20,xmax=44,ymin=-0.5,ymax=7,grid=major,grid style={gray!25},legend pos=south east,
  legend style={font=\scriptsize},table/col sep=comma]
\addplot[omThm,thick] table[x=day,y=early02]{figdata/strategies-1/10-index-rebalancing/path.csv};
\addplot[omDef,very thick] table[x=day,y=early05]{figdata/strategies-1/10-index-rebalancing/path.csv};
\addplot[black!50,thick,dashed] table[x=day,y=early08]{figdata/strategies-1/10-index-rebalancing/path.csv};
\addplot[black!40,dotted] coordinates {(0,-0.5) (0,7)};
\addplot[black!40,dotted] coordinates {(24,-0.5) (24,7)};
\legend{early share 0.2,early share 0.5,early share 0.8}
\end{axis}
\end{tikzpicture}

\omcaption{Additions to the synthetic index, 15\% of shares tracked: mean market-adjusted return with the planted index-fund effect, from twenty trading days
before the announcement (day 0) through the effective close (day 24) and the twenty days after, for three shares of the effect taken before the announcement.
Data: \texttt{s1\_index.path\_for\_figure}.}\label{fig:s1:index-rebalancing:path}
\end{omfigure}

Three trades share the event. \emph{Predicting} buys additions (and sells deletions) twenty days before the announcement and holds to it: with a perfect
prediction it earns 3.75\% per event at an early share of one half. \emph{Trading the announcement} buys at the announcement's close and sells at the effective
close to the \omterm{def:s1:index-rebalancing:fund}{index funds}: 2.40\% per event after costs (standard error 0.4\%). \emph{Providing liquidity at the effective close}, selling to the \omterm{def:s1:index-rebalancing:fund}{index funds} and
buying back over the next twenty days, earns the reversal: 2.65\%. The additions keep a small drift of their own, since stocks that have grown into an index
are recent winners and the synthetic market plants momentum; Chen, Noronha and Singal found the real asymmetry the other way round, with a permanent rise
for additions and no permanent fall for deletions, which they attributed partly to investor awareness.

\section{How the trade decayed}

\begin{center}
\small
\begin{tabular}{@{}lccc@{}}
\toprule
return from the announcement to the effective close & 5\% of shares tracked & 15\% & 30\%\\
\midrule
early share 0.2 & 2.2\% & 4.0\% & 5.7\%\\
early share 0.5 & 1.3\% & 2.4\% & 3.5\%\\
early share 0.8 & 0.4\% & 0.8\% & 1.3\%\\
\midrule
average push & 3.1\% & 5.3\% & 7.5\%\\
\bottomrule
\end{tabular}
\end{center}

The event return rises with the demand and falls with the share taken early. More indexed money makes the push larger: from 3.1\% at 5\% of shares tracked to
7.5\% at 30\%. More arbitrage capital predicting the change moves it earlier, out of reach of anyone who waits for the announcement: at an early share of 0.8
the announcement trade keeps a fifth of what it keeps at 0.2. The predictors' own return rises to 5.3\% per event, but only with perfect foresight; twenty
days out a predictor is wrong on one addition in seven and misses more than a quarter. This is the pattern the public record shows for the S\&P 500: an
effect first documented as a permanent shift in demand, and later as a temporary one that reverts, in a market where the event is known to all. The trade
did not disappear; it moved earlier, to those who predict, and later, to those who provide liquidity at the close.

\section{Strategy files}

\begin{strategyfile}{Predicted additions}\label{strat:s1:index-rebalancing:additions}
\sfield{Who pays you, and why} \omterm{def:s1:index-rebalancing:fund}{Index funds} that must buy at the effective close, whatever the price.
\sfield{Instruments and venues} Stocks near the index cut-off; the closing auction on the effective date.
\sfield{Signal} The probability of joining, from the index rules applied to projected characteristics.
\sfield{Sizing and execution} Buy likely additions before the \omterm{def:s1:index-rebalancing:fund}{rank day} or the announcement; sell into the \omterm{def:s1:index-rebalancing:fund}{index funds}' demand at the effective close.
\sfield{Costs} Trading in stocks around the cut-off, often less liquid; wrong predictions.
\sfield{How it dies} Competition moving the price earlier; \omterm{def:s1:index-rebalancing:fund}{index funds} that trade over several days.
\sfield{Horizon, capacity, infrastructure} Weeks; capacity limited by the event's size; an index-rule engine and capitalisation data.
\sfield{Backtest honestly} Rules as they were at each date; predictions from data before the \omterm{def:s1:index-rebalancing:fund}{rank day}; the wrong predictions counted.
\sfield{Sources} Shleifer (1986); this chapter's simulation (3.75\% per event with perfect prediction at an early share of one half).
\end{strategyfile}

\begin{strategyfile}{Predicted deletions}\label{strat:s1:index-rebalancing:deletions}
\sfield{Who pays you, and why} \omterm{def:s1:index-rebalancing:fund}{Index funds} that must sell at the effective close.
\sfield{Instruments and venues} Members near the exit threshold.
\sfield{Signal} The probability of leaving.
\sfield{Sizing and execution} Short likely deletions (borrow permitting) and buy them back from the \omterm{def:s1:index-rebalancing:fund}{index funds}' selling at the close.
\sfield{Costs} Borrow on falling, often small, stocks.
\sfield{How it dies} As for additions; deleted stocks may not stay down.
\sfield{Horizon, capacity, infrastructure} Weeks; borrow availability.
\sfield{Backtest honestly} Borrow fees and availability; delisting of the weakest members.
\sfield{Sources} Chen, Noronha and Singal (2004): no permanent decline for deletions.
\end{strategyfile}

\begin{strategyfile}{Effective-date close liquidity provision}\label{strat:s1:index-rebalancing:liquidity}
\sfield{Who pays you, and why} \omterm{def:s1:index-rebalancing:fund}{Index funds} that trade at one close regardless of price.
\sfield{Instruments and venues} The closing auction (Book~1, chapter~13) on the effective date.
\sfield{Signal} The published imbalance of the closing auction and the known index demand.
\sfield{Sizing and execution} Sell additions and buy deletions in the auction; unwind over the following weeks.
\sfield{Costs} Holding risk over the reversal.
\sfield{How it dies} When \omterm{def:s1:index-rebalancing:fund}{index funds} spread their trading, or too much capital meets them at the close.
\sfield{Horizon, capacity, infrastructure} Days to weeks; auction imbalance data.
\sfield{Backtest honestly} Auction prices, not the day's average; capacity against the auction's size.
\sfield{Sources} This chapter's simulation (2.65\% per event when half the push reverses); Patel and Welch (2017) on reversion in the 2000s.
\end{strategyfile}

\begin{strategyfile}{Post-event reversal}\label{strat:s1:index-rebalancing:reversal}
\sfield{Who pays you, and why} As for liquidity provision: the part of the push that was pressure, not information.
\sfield{Instruments and venues} Recently added and deleted stocks.
\sfield{Signal} The size of the push relative to the stock's volume.
\sfield{Sizing and execution} Short additions and buy deletions after the effective date, hedged.
\sfield{Costs} Borrow on additions; the momentum of recent winners.
\sfield{How it dies} If index membership carries information (awareness, liquidity), the push does not revert.
\sfield{Horizon, capacity, infrastructure} Weeks to months.
\sfield{Backtest honestly} Separate additions from deletions; control for momentum.
\sfield{Sources} Patel and Welch (2017); Chen, Noronha and Singal (2004).
\end{strategyfile}

\section{Tutorial: everyone predicts it}\label{tut:s1:index-rebalancing:tut}

\begin{tutorial}
\textbf{Goal.} Write a banded index rule, predict its changes, and trade the event at three moments. \textbf{End state:} the two tables and
\cref{fig:s1:index-rebalancing:path}.
\begin{tutsteps}
\item \textbf{The rule and its prediction}: ranks, the band, and forward simulation of capitalisations.
\omcode{code/firm/indexevent/firm_indexevent.py}{27}{48}{The banded reconstitution rule and its prediction.}\label{lst:s1:index-rebalancing:rule}
\item \textbf{The event}: the planted path from prediction to reversal.
\omcode{code/firm/indexevent/firm_indexevent.py}{55}{63}{The event's price path.}\label{lst:s1:index-rebalancing:path}
\item \textbf{Run} \texttt{accuracy} at 0, 20 and 60 days, \texttt{trades} over the grid, and \texttt{fig\_index.py}.
\end{tutsteps}
\textbf{What to change next.} Trade only predictions above a probability of 0.8 and count the wrong ones; let \omterm{def:s1:index-rebalancing:fund}{index funds} trade over three days; add a
semi-annual and an annual calendar and compare the churn.
\end{tutorial}

\section{Build: index events}\label{bld:s1:index-rebalancing:indexevent}

\begin{build}
\bfield{Purpose} Index rules, their prediction, and the price path of index-fund demand, to study index events with a known truth.
\bfield{Interface} \texttt{reconstitute(cap, members, size, band)}, \texttt{predict(cap, vol, members, size, band, days, sims, rng)},
\texttt{event\_push(q, adv, sigma, eta)}, \texttt{event\_path(push, early, revert, pre, run, post)}.
\bfield{Rules} Membership from rank-day data only; predictions from earlier data; the event path added to market-adjusted returns.
\bfield{Acceptance tests} \texttt{code/firm/indexevent/tests/}: the band rule by hand; certain predictions far from the cut; the push and the path by hand.
\bfield{Stretch} Float adjustment and liquidity screens; committee-style discretionary changes; \omterm{def:s1:index-rebalancing:fund}{index funds} that trade over several days.
\end{build}

\begin{omsources}
\item A.~Shleifer, ``Do demand curves for stocks slope down?'', \emph{Journal of Finance} 41(3), 1986.
\item L.~Harris and E.~Gurel, ``Price and volume effects associated with changes in the S\&P 500 list'', \emph{Journal of Finance} 41(4), 1986.
\item H.~Chen, G.~Noronha and V.~Singal, ``The price response to S\&P 500 index additions and deletions'', \emph{Journal of Finance} 59(4), 2004.
\item N.~Patel and I.~Welch, ``Extended stock returns in response to S\&P 500 index changes'', \emph{Review of Asset Pricing Studies} 7(2), 2017.
\item FTSE Russell (LSEG), Russell US indexes reconstitution announcements, 2026.
\end{omsources}

\section{Exercises}

\begin{exercise}[$\star$]\label{exo:s1:index-rebalancing:1}
An index of 300 stocks has a band of 15 ranks. At what rank does a non-member join, and at what rank does a member leave?
\end{exercise}

\begin{exercise}[$\star$]\label{exo:s1:index-rebalancing:2}
\omterm{def:s1:index-rebalancing:fund}{Index funds} must buy 15\% of a stock's shares; it trades 0.6\% of its shares a day with a daily volatility of 2\%. What push does the square-root law give?
\end{exercise}

\begin{exercise}[$\star$]\label{exo:s1:index-rebalancing:3}
With half of that push taken before the announcement and half of it reversing after the effective date, what do the announcement trade and the liquidity
provider earn before costs?
\end{exercise}

\begin{exercise}[$\star\star$]\label{exo:s1:index-rebalancing:4}
Why does a band reduce turnover, and what does it do to the predictability of changes?
\end{exercise}

\begin{exercise}[$\star\star$]\label{exo:s1:index-rebalancing:5}
Explain why the event return from the announcement falls as arbitrage capital grows, while the event itself does not shrink.
\end{exercise}

\begin{exercise}[$\star\star$]\label{exo:s1:index-rebalancing:6}
Predictions sixty days before the \omterm{def:s1:index-rebalancing:fund}{rank day} are right 76\% of the time. What does that do to a predictor's return, compared with a perfect predictor?
\end{exercise}

\begin{exercise}[$\star\star\star$]\label{exo:s1:index-rebalancing:7}
\emph{Coding.} Run \texttt{trades(0.30, 0.2)} and \texttt{trades(0.05, 0.8)}. Explain the difference in the announcement trade's return.
\end{exercise}

\begin{exercise}[$\star\star\star$]\label{exo:s1:index-rebalancing:8}
\emph{Find the flaw.} ``Stocks added to the index rose 7\% from the announcement to the effective date in our 1990s sample; we will buy all announced
additions.''
\end{exercise}

\section{Problem: Everyone Predicts It}

\begin{problem}[{Weekend problem --- a known event, traded}]\label{pb:s1:index-rebalancing:1}
The synthetic index and the public record.

\medskip\noindent\textbf{Part I --- The rules.}
\begin{enumerate}
\item Define an \omterm{def:s1:index-rebalancing:fund}{index fund}, a \omterm{def:s1:index-rebalancing:fund}{rank day} and \omterm{def:s1:index-rebalancing:prediction}{membership prediction}.
\item Describe the Russell calendar in 2026.
\item Describe the synthetic index's rule and calendar, and the number of changes.
\item What does a band do?
\end{enumerate}

\medskip\noindent\textbf{Part II --- Prediction.}
\begin{enumerate}[resume]
\item How are changes predicted before the \omterm{def:s1:index-rebalancing:fund}{rank day}?
\item Give the precision and recall at 0, 20 and 60 days.
\item Why does prediction get worse with horizon?
\item What is the trade-off between early and certain?
\end{enumerate}

\medskip\noindent\textbf{Part III --- The event.}
\begin{enumerate}[resume]
\item Describe the planted event path.
\item Give the returns of predicting, trading the announcement and providing liquidity.
\item What did Shleifer and Chen, Noronha and Singal find?
\item Why do synthetic additions keep a small drift?
\end{enumerate}

\medskip\noindent\textbf{Part IV --- The verdict.}
\begin{enumerate}[resume]
\item State the \emph{named result}: the prediction accuracy of membership changes and the event return as index-tracking demand grows relative to
arbitrage capital.
\item What did Patel and Welch find?
\item Where did the trade go?
\item Which strategy file would you run with small capital, and which with large?
\item What would you monitor on the effective date?
\item How do semi-annual reconstitutions change the trade?
\item What makes an index event different from a flow event of chapter~9?
\item In one sentence: what is the index effect today?
\end{enumerate}
\end{problem}

\section{Interview questions}

\begin{interviewq}[$\star$ \iqroles{researcher, trader}]\label{iq:s1:index-rebalancing:1}
Why does a stock's price move when it is added to an index?
\end{interviewq}

\begin{interviewq}[$\star\star$ \iqroles{trader}]\label{iq:s1:index-rebalancing:2}
How would you trade the Russell reconstitution?
\end{interviewq}

\begin{interviewq}[$\star\star$ \iqroles{researcher, developer}]\label{iq:s1:index-rebalancing:3}
Build a predictor of index additions a month before the \omterm{def:s1:index-rebalancing:fund}{rank day}. What inputs and what uncertainty?
\end{interviewq}

\begin{interviewq}[$\star\star$ \iqroles{risk}]\label{iq:s1:index-rebalancing:4}
What can go wrong with a book of predicted index additions?
\end{interviewq}

\begin{interviewq}[$\star\star$ \iqroles{trader}]\label{iq:s1:index-rebalancing:5}
The closing auction on reconstitution day shows a large buy imbalance in a stock. What do you do?
\end{interviewq}

\begin{interviewq}[$\star\star\star$ \iqroles{researcher}]\label{iq:s1:index-rebalancing:6}
Index demand $D$ meets arbitrage capital $A$ that buys before the announcement. In a model where the price push is proportional to the net demand the
market must absorb at each stage, how does the return after the announcement depend on $A/D$?
\end{interviewq}
